\chapter{Langlands--Shahidi Normalisation of the Eisenstein Koszul Density}
\label{v4-arith:w6-c:LS-correction}

\emph{The bounded-interval reduction, construction~C. Chapter~\ref{v4-arith:w5-c:full-vprim-P3}
(the restricted-class comparison) reduced Koszul--Petersson compatibility on the
discrete-spectrum core by domain restriction: the Eisenstein
continuous sector was placed outside the P3 identity, and the Li
trace was separated from it under the H-P\(^{\mathrm{Ar}}\) spectral
realisation. The Langlands--Shahidi alternative was sketched as
the second normalisation of \S\ref{v4-arith:w5-c:LS}. The exact
conditional normalisation problem is this: multiplication by a
Langlands--Shahidi scalar removes the unitary phase only after a
choice of positive Plancherel measure has been fixed; the boundary
Euler product on \(\mathrm{Re}(s)=1/2\) is meromorphic, not absolutely
convergent. Thus the Langlands--Shahidi normalisation is a conditional construction, not
an unconditional closure of the Eisenstein residual.}

\section{Setup and the LS Normalisation}
\label{v4-arith:w6-c:LS:setup}

\subsection{The Residual Mismatch}
\label{v4-arith:w6-c:LS:residual}

The explicit mismatch of
\Cref{v4-arith:w5-c:prop:mismatch}
(Chapter~\ref{v4-arith:w5-c:full-vprim-P3})
is as follows. On the Eisenstein continuous spectrum
\(L^{2}_{\mathrm{cont}}([GL_{2}])\) parametrised by
\(\mathrm{Re}(s)=1/2\), the Petersson-side Maass--Selberg normalisation and
the Koszul-side bar--cobar coefficient are recorded in the
unnormalised spherical convention by
\begin{align}
\label{v4-arith:w6-c:eq:Pet-density}
\operatorname{lead}^{\mathrm{Pet}}_{\mathrm{cont}}(s)
&\;=\;\bigl(1+|M(s)|^{2}\bigr)|\xi(2s)|^{-2}\qquad(\mathrm{Re}(s)=1/2),\\
\label{v4-arith:w6-c:eq:Koszul-density}
\operatorname{lead}^{\mathrm{Koszul}}_{\mathrm{cont}}(s)
&\;=\;\frac{\xi(2s)}{\xi(2s-1)}\cdot|\xi(2s)|^{-2}\qquad(\mathrm{Re}(s)=1/2).
\end{align}
The ratio \(\xi(2s)/\xi(2s-1)\) has modulus~\(1\) on the critical line
(lemma below). Thus the mismatch is a normalisation and phase problem
between the Maass--Selberg scattering coefficient and the Koszul
coefficient; it is not removed by treating the diagonal as an ordinary
\(L^{2}\)-density.

\begin{lemma}[unit modulus of \(\xi(2s)/\xi(2s-1)\) on the critical line]
\label{v4-arith:w6-c:lem:unit-modulus}
\ClaimStatusProvedHere. For \(s=1/2+it\), away from the Tate boundary
and common zeros of the two factors,
\begin{equation}
\label{v4-arith:w6-c:eq:unit-modulus}
\bigl|\xi(2s)/\xi(2s-1)\bigr|\;=\;1.
\end{equation}
\end{lemma}

\begin{proof}
For \(s=1/2+it\), write \(2s=1+2it\) and \(2s-1=2it\). The completed
zeta \(\xi\) satisfies \(\xi(w)=\xi(1-w)\) (functional equation;
Titchmarsh 1986~\S2.1), giving \(\xi(1+2it)=\xi(-2it)\). Complex
conjugation gives \(\overline{\xi(1+2it)}=\xi(1-2it)\) (Schwarz
reflection: \(\xi\) is real on the real axis, hence
\(\overline{\xi(w)}=\xi(\overline{w})\)). Combining,
\(\xi(1+2it)\cdot\overline{\xi(1+2it)}=\xi(1+2it)\xi(1-2it)=\xi(1+2it)\xi(2it)\)
by the functional equation. Similarly
\(\xi(2it)\cdot\overline{\xi(2it)}=\xi(2it)\cdot\xi(\overline{2it})=\xi(2it)\cdot\xi(-2it)=\xi(2it)\cdot\xi(1+2it)\)
by FE. Therefore \(|\xi(2s)|^{2}=|\xi(2s-1)|^{2}\), giving
(\ref{v4-arith:w6-c:eq:unit-modulus}).
\end{proof}

\subsection{The Langlands--Shahidi Normalisation Factor}
\label{v4-arith:w6-c:LS:factor-def}

We follow Shahidi 2010 \emph{Eisenstein Series and Automorphic
\(L\)-Functions} Chapter~2 and Moeglin--Waldspurger 1995 \S IV.3 for
the intertwining operator normalisation.

\begin{definition}[Langlands--Shahidi normalisation factor on \(GL_{2}\) Eisenstein]
\label{v4-arith:w6-c:def:LS-factor}
The \emph{Langlands--Shahidi normalisation factor} on the
Eisenstein summand \(\mathrm{Ind}_{B}^{G}(\chi_{s})\) of \(GL_{2}\) at
spectral parameter \(s\) with \(\mathrm{Re}(s)=1/2\) is
\begin{equation}
\label{v4-arith:w6-c:eq:LS-factor}
N^{\mathrm{LS}}(s)\;:=\;\frac{\xi(2s-1)}{\xi(2s)}.
\end{equation}
Its inverse is
\(N^{\mathrm{LS}}(s)^{-1}=\xi(2s)/\xi(2s-1)\).
\end{definition}

\begin{remark}[origin and canonicity of the factor]
\label{v4-arith:w6-c:rem:LS-origin}
The factor \(N^{\mathrm{LS}}(s)\) is the global Langlands--Shahidi
normalisation of the standard intertwining operator
\(M(s)\colon\mathrm{Ind}_{B}^{G}(\chi_{s})\to\mathrm{Ind}_{B}^{G}(\chi_{s}^{w})\),
with three essential properties.
\begin{enumerate}
\item[(LS1)] \emph{Local factor decomposition.} For each place~\(v\),
\(N^{\mathrm{LS}}(s)\) is a product of local Langlands factors
\(N^{\mathrm{LS}}_{v}(s)\) as specified by Casselman--Shalika 1980
\emph{Compositio Math.}~41 at unramified finite~\(v\) and by Shahidi
2010~\S2.2 at ramified and archimedean~\(v\). Restricted tensor
product factorisation: \(N^{\mathrm{LS}}(s)=\prod_{v}N^{\mathrm{LS}}_{v}(s)\).
\item[(LS2)] \emph{Functional equation.}
\(N^{\mathrm{LS}}(s)\cdot N^{\mathrm{LS}}(1-s)=1\)
on the critical line. Explicitly:
\[
N^{\mathrm{LS}}(s)\cdot N^{\mathrm{LS}}(1-s)
=\frac{\xi(2s-1)}{\xi(2s)}\cdot\frac{\xi(1-2s)}{\xi(2-2s)}
=\frac{\xi(2s-1)\cdot\xi(2s)}{\xi(2s)\cdot\xi(2s-1)}=1,
\]
using the functional equation \(\xi(1-2s)=\xi(2s)\) and
\(\xi(2-2s)=\xi(2s-1)\) (Titchmarsh 1986~\S2.1). The Shahidi
global normalisation satisfies \(n(s)\cdot n(1-s)=1\) (Shahidi
2010~Thm~3.5.1). No Maass--Selberg measure scalar is part of
\(N^{\mathrm{LS}}\).
\item[(LS3)] \emph{Unitarity of the normalised intertwiner.} After
division by \(N^{\mathrm{LS}}(s)\), the normalised intertwiner
\(\widetilde{M}(s)=M(s)/N^{\mathrm{LS}}(s)\) is unitary on the
principal-series representation at each critical parameter
(Moeglin--Waldspurger~\S IV.3.10).
\end{enumerate}
\end{remark}

\section{The Redefined Koszul Pairing on the Eisenstein Summand}
\label{v4-arith:w6-c:LS:redefinition}

\begin{definition}[LS-corrected Koszul pairing on the Eisenstein summand]
\label{v4-arith:w6-c:def:LS-Koszul}
Let \(\phi_{s_{1}},\phi_{s_{2}}\in\mathrm{Ind}_{B}^{G}(\chi_{s})\)
be spherical vectors at Eisenstein parameters \(s_{1},s_{2}\) with
\(\mathrm{Re}(s_{i})=1/2\). The \emph{Langlands--Shahidi-corrected
Koszul pairing} is
\begin{equation}
\label{v4-arith:w6-c:eq:LS-Koszul-def}
\bigl\langle\phi_{s_{1}},\sigma^{\mathrm{BC}}(\phi_{s_{2}})\bigr\rangle^{\mathrm{Koszul},\mathrm{LS}}
\;:=\;N^{\mathrm{LS}}(\tfrac{1}{2}(s_{1}+s_{2}))^{-1}\cdot\bigl\langle\phi_{s_{1}},\sigma^{\mathrm{BC}}(\phi_{s_{2}})\bigr\rangle^{\mathrm{Koszul}}.
\end{equation}
On the diagonal \(s_{1}=s_{2}=s\) this reduces to
\begin{equation}
\label{v4-arith:w6-c:eq:LS-Koszul-diagonal}
\bigl\langle\phi_{s},\sigma^{\mathrm{BC}}(\phi_{s})\bigr\rangle^{\mathrm{Koszul},\mathrm{LS}}
\;=\;\frac{\xi(2s)}{\xi(2s-1)}\cdot\bigl\langle\phi_{s},\sigma^{\mathrm{BC}}(\phi_{s})\bigr\rangle^{\mathrm{Koszul}}.
\end{equation}
\end{definition}

\begin{theorem}[conditional LS Plancherel normalisation match on the continuum]
\label{v4-arith:w6-c:thm:density-match}
\ClaimStatusProvedHereConditional \emph{(conditional on the
bar--cobar Eisenstein pairing being transported to the positive
Langlands--Shahidi Plancherel measure)}. For spherical vectors
\(\phi_{s}\in\mathrm{Ind}_{B}^{G}(\chi_{s})\) with \(\mathrm{Re}(s)=1/2\),
\begin{equation}
\label{v4-arith:w6-c:eq:density-match}
\mathsf{Sh}_{\mathrm{LS}}\!\left(\operatorname{lead}^{\mathrm{Koszul},\mathrm{LS}}_{\mathrm{cont}}(s)\right)
\;=\;\operatorname{lead}^{\mathrm{Pet}}_{\mathrm{cont}}(s).
\end{equation}
Equivalently,
\begin{equation}
\label{v4-arith:w6-c:eq:density-match-equiv}
\mathsf{Sh}_{\mathrm{LS}}\!\left(\mathrm{Koszul\ continuum}\right)
\;=\;\mathrm{Petersson\ continuum}\qquad(\mathrm{Re}(s)=1/2),
\end{equation}
after this measure identification. Here \(\mathsf{Sh}_{\mathrm{LS}}\)
denotes the Shahidi Plancherel transport. No pointwise equality of raw
scalar densities is asserted before applying \(\mathsf{Sh}_{\mathrm{LS}}\).
\end{theorem}

\begin{proof}
The Koszul pairing of
\Cref{v4-arith:kms-gns:lem:Koszul-continued} gives, on
\(s_{1}=s_{2}=s\) after the common pole at \(2\mathrm{Re}(s)-1=0\) is
regularised,
\begin{equation*}
\bigl\langle\phi_{s},\sigma^{\mathrm{BC}}(\phi_{s})\bigr\rangle^{\mathrm{Koszul}}
\;=\;\frac{\xi(2s)}{\xi(2s-1)}\cdot\bigl\langle\phi_{s},\phi_{s}\bigr\rangle^{\mathrm{loc}}.
\end{equation*}
The ratio \(\xi(2s)/\xi(2s-1)\) has modulus one on the critical line,
so it is a scattering phase, not a positive density. Therefore the
desired comparison cannot be a consequence of multiplying scalar
densities. It is precisely the conditional assertion that the
bar--cobar Eisenstein quadratic form is transported to the Shahidi
positive Plancherel measure. With that transport supplied,
\eqref{v4-arith:w6-c:eq:density-match} is the definition of the
matched normalisation.
\end{proof}

\begin{remark}[why no scalar square proves the match]
\label{v4-arith:w6-c:rem:squared-modulus}
The diagonal pairing is only a distributional Maass--Selberg object
before truncation is regularised. Since
\(\xi(2s)/\xi(2s-1)\) has unit modulus on the critical line, taking
absolute values removes the phase but does not identify the underlying
Plancherel normalisation. The non-formal point is exactly the
identification of the bar--cobar Eisenstein quadratic form with the
Shahidi positive Plancherel measure.
\end{remark}

\begin{corollary}[conditional normalisation match via LS]
\label{v4-arith:w6-c:cor:unconditional}
\ClaimStatusProvedHereConditional \emph{(same hypothesis as
\Cref{v4-arith:w6-c:thm:density-match})}. Under the LS redefinition
\Cref{v4-arith:w6-c:def:LS-Koszul}, the Koszul and Petersson
normalisations on the Eisenstein continuous spectrum of
\(L^{2}([GL_{2}])\) coincide after the positive Plancherel-measure
identification.
\end{corollary}

\begin{proof}
\Cref{v4-arith:w6-c:thm:density-match} is the stated
measure-identification hypothesis written in the LS normalisation. No
Ramanujan input is used on the continuum; the continuum is
Eisenstein-induced and tempered on the unitary axis by construction.
\end{proof}

\section{Compatibility Checks}
\label{v4-arith:w6-c:LS:compat}

Three compatibility properties of the LS-corrected pairing are
verified: the arithmetic Heisenberg Koszul self-duality
\(\kappa^{\mathrm{Ar}}+(\kappa^{!})^{\mathrm{Ar}}=0\) on the
Eisenstein summand (check~(i)); the Tate restricted-tensor
factorisation of \(N^{\mathrm{LS}}\) across places (check~(ii)); and
the compatibility with the discrete-core result of
Chapter~\ref{v4-arith:w5-c:full-vprim-P3} (check~(iii)).

\subsection{Compatibility Check (i): Heisenberg Self-Duality}
\label{v4-arith:w6-c:LS:compat-koszul}

\begin{proposition}[LS-corrected Koszul preserves Heisenberg self-duality on the Eisenstein summand]
\label{v4-arith:w6-c:prop:koszul-selfduality}
\ClaimStatusProvedHereConditional \emph{(conditional on the LS
scalar being admitted as part of the arithmetic Koszul pairing on
the Eisenstein summand)}. The LS-corrected Koszul pairing
\(\langle\,\cdot\,,\,\cdot\,\rangle^{\mathrm{Koszul},\mathrm{LS}}\)
of \Cref{v4-arith:w6-c:def:LS-Koszul} satisfies the arithmetic
Heisenberg Koszul self-duality
\begin{equation}
\label{v4-arith:w6-c:eq:koszul-selfduality}
\kappa^{\mathrm{Ar}}_{\mathsf{G}}+(\kappa^{!})^{\mathrm{Ar}}_{\mathsf{G}}\;=\;0
\end{equation}
on the Eisenstein summand, where \(\mathsf{G}=GL_{2}\) and
\(\kappa^{\mathrm{Ar}},(\kappa^{!})^{\mathrm{Ar}}\) are the
arithmetic Heisenberg and dual-Heisenberg
Koszul obstructions of
\Cref{v4-arith:thm:P1-resolution}.
\end{proposition}

\begin{proof}
The self-duality identity \(\kappa^{\mathrm{Ar}}+(\kappa^{!})^{\mathrm{Ar}}=0\)
is equivalent to the completed-zeta functional equation
\(\xi(s)=\xi(1-s)\) (\Cref{v4-arith:w5-c:full-vprim-P3} and
\Cref{v4-arith:cl:eq:G-row}). The LS correction
\(N^{\mathrm{LS}}(s)\) is a ratio of completed zetas:
\(N^{\mathrm{LS}}(s)=\xi(2s-1)/\xi(2s)\). Its own
functional equation (LS2):
\begin{equation*}
N^{\mathrm{LS}}(s)\cdot N^{\mathrm{LS}}(1-s)\;=\;\frac{\xi(2s-1)}{\xi(2s)}\cdot\frac{\xi(1-2s)}{\xi(2-2s)}\;=\;\frac{\xi(2s-1)\xi(1-2s)}{\xi(2s)\xi(2-2s)}\;=\;\frac{\xi(2s-1)\xi(2s)}{\xi(2s)\xi(2s-1)}\;=\;1
\end{equation*}
using \(\xi(1-2s)=\xi(2s)\) and \(\xi(2-2s)=\xi(2s-1)\)
(Titchmarsh 1986~\S2.1). The self-duality
\(\kappa^{\mathrm{Ar}}+(\kappa^{!})^{\mathrm{Ar}}=0\) is an
obstruction-class identity, equivalent to the completed-zeta
functional equation \(\xi(s)=\xi(1-s)\)
(\Cref{v4-arith:cl:eq:G-row}: the sign
\(\kappa^{\mathrm{Ar}}_{\mathsf{G}}+(\kappa^{!})^{\mathrm{Ar}}_{\mathsf{G}}=0\)
is precisely the G-row sign of the arithmetic Heisenberg duality,
independent of rescaling of the pairing). The LS correction acts on the pairing level, not on the
obstruction-class level: the Koszul obstructions
\(\kappa^{\mathrm{Ar}},(\kappa^{!})^{\mathrm{Ar}}\) are defined
as obstruction classes in the arithmetic bar--cobar complex
(\Cref{v4-arith:thm:P1-resolution}), independent of the choice of
inner product on \(\mathcal{V}^{\mathrm{prim}}_{\mathrm{cont}}\),
and hence invariant under multiplication of the Koszul pairing
by \(N^{\mathrm{LS}}(s)^{-1}\). The functional equation
\(\xi(s)=\xi(1-s)\) (which holds unconditionally) therefore gives
\(\kappa^{\mathrm{Ar}}+(\kappa^{!})^{\mathrm{Ar}}=0\) on the
Eisenstein summand after the LS correction, as required.
\end{proof}

\begin{remark}[no modification to the discrete sector]
\label{v4-arith:w6-c:rem:discrete-untouched}
The LS correction factor is by construction supported on the
Eisenstein continuous summand: \(N^{\mathrm{LS}}(s)^{-1}=1\) at
discrete-spectrum points (pole at the cuspidal locus gives
\(N^{\mathrm{LS}}\) a compensating zero, and the discrete spectrum
is outside the domain of the continuous normalisation).
The discrete-spectrum Koszul pairing of
Chapter~\ref{v4-arith:w5-c:full-vprim-P3} is therefore unchanged.
\end{remark}

\subsection{Compatibility Check (ii): Tate Factorisation}
\label{v4-arith:w6-c:LS:compat-tate}

\begin{proposition}[LS factorises over places as a Tate restricted tensor product]
\label{v4-arith:w6-c:prop:tate-factorisation}
\ClaimStatusProvedHereConditional \emph{(as a meromorphic
Langlands--Shahidi factorisation, not as an absolutely convergent
Euler product on the critical line)}. The global Langlands--Shahidi factor
\(N^{\mathrm{LS}}(s)\) of \Cref{v4-arith:w6-c:def:LS-factor} admits
a place-by-place decomposition
\begin{equation}
\label{v4-arith:w6-c:eq:LS-local-factors}
N^{\mathrm{LS}}(s)\;=\;\prod_{v\in|\overline{\mathrm{Spec}\,\mathbb{Z}}|}N^{\mathrm{LS}}_{v}(s),
\end{equation}
with local factors
\begin{equation}
\label{v4-arith:w6-c:eq:LS-local-factor-v}
N^{\mathrm{LS}}_{v}(s)\;=\;\begin{cases}
\zeta_{p}(2s-1)/\zeta_{p}(2s)\qquad&(v=p\text{ finite, unramified})\\
\zeta_{p}(2s-1)/\zeta_{p}(2s)\cdot\varepsilon_{p}(s,\pi_{p},\psi_{p})\qquad&(v=p\text{ finite, ramified})\\
\Gamma_{\mathbb{R}}(2s-1)/\Gamma_{\mathbb{R}}(2s)\qquad&(v=\infty),
\end{cases}
\end{equation}
where \(\zeta_{p}(s)=(1-p^{-s})^{-1}\) is the Euler factor,
\(\Gamma_{\mathbb{R}}(s)=\pi^{-s/2}\Gamma(s/2)\) the archimedean
Tate factor, and \(\varepsilon_{p}\) the local
Jacquet--Piatetski-Shapiro--Shalika epsilon factor at the
principal-series induction. The product is interpreted by
meromorphic continuation of the completed \(L\)-factor.
\end{proposition}

\begin{proof}
The completed zeta factorises
\(\Lambda_{\zeta}(s)=\Gamma_{\mathbb{R}}(s)\prod_{p}\zeta_{p}(s)\)
(Tate 1950~\S4.2). Applying the factorisation to the ratio
\(\Lambda_{\zeta}(2s-1)/\Lambda_{\zeta}(2s)\):
\begin{equation*}
\frac{\Lambda_{\zeta}(2s-1)}{\Lambda_{\zeta}(2s)}\;=\;\frac{\Gamma_{\mathbb{R}}(2s-1)}{\Gamma_{\mathbb{R}}(2s)}\cdot\prod_{p}\frac{\zeta_{p}(2s-1)}{\zeta_{p}(2s)}.
\end{equation*}
Unramified local factors are \(\zeta_{p}(2s-1)/\zeta_{p}(2s)\);
ramified local factors pick up an additional local epsilon factor
\(\varepsilon_{p}\) from the Jacquet--Piatetski-Shapiro--Shalika 1983
local-functional-equation correction, which has absolute
value~\(1\) and does not affect the Plancherel normalisation match
(\Cref{v4-arith:w6-c:thm:density-match}). The restricted tensor
product \(\prod_{v}N^{\mathrm{LS}}_{v}(s)\) is not an absolutely
convergent Euler product on \(\mathrm{Re}(s)=1/2\); it is the
boundary value of the meromorphic completed ratio, giving
(\ref{v4-arith:w6-c:eq:LS-local-factors}).
\end{proof}

\begin{remark}[Casselman--Shalika at unramified finite places]
\label{v4-arith:w6-c:rem:casselman-shalika}
The local factor at an unramified finite prime~\(p\),
\(N^{\mathrm{LS}}_{p}(s)=\zeta_{p}(2s-1)/\zeta_{p}(2s)=(1-p^{-2s})/(1-p^{1-2s})\),
is precisely the Casselman--Shalika 1980 formula for the
unramified spherical Whittaker function at principal-series
parameter~\(s\) (Casselman--Shalika \emph{Compositio Math.}~41
Thm~5.4 with \(w\) the non-trivial Weyl element). This identifies
\(N^{\mathrm{LS}}_{p}\) as the local intertwining operator
normalisation, canonically and without convention choice.
\end{remark}

\subsection{Compatibility Check (iii): Discrete-Core Fixation}
\label{v4-arith:w6-c:LS:compat-discrete}

\begin{proposition}[LS correction preserves the discrete-core identity]
\label{v4-arith:w6-c:prop:discrete-fixation}
\ClaimStatusProvedHereConditional \emph{(with the hypotheses of the
discrete-core identity)}. The discrete-spectrum Koszul--Petersson
identity of
\Cref{v4-arith:w5-c:thm:discrete-core-identity}
is preserved under the substitution
\(\langle\,\cdot\,,\,\cdot\,\rangle^{\mathrm{Koszul}}\leadsto
\langle\,\cdot\,,\,\cdot\,\rangle^{\mathrm{Koszul},\mathrm{LS}}\)
extended to the Eisenstein summand:
\begin{equation}
\label{v4-arith:w6-c:eq:discrete-preserved}
\langle v,w\rangle^{\mathrm{Pet}}\;=\;\langle v,\sigma^{\mathrm{BC}}(w)\rangle^{\mathrm{Koszul},\mathrm{LS}}\qquad\forall v,w\in\mathcal{V}^{\mathrm{prim}}_{\mathrm{disc}}.
\end{equation}
\end{proposition}

\begin{proof}
\(\mathcal{V}^{\mathrm{prim}}_{\mathrm{disc}}\) is orthogonal to the
Eisenstein continuous summand in the Langlands decomposition
(Langlands 1966~\S6; Moeglin--Waldspurger 1995 Ch.~II). The
substitution
\(\langle\,\cdot\,,\,\cdot\,\rangle^{\mathrm{Koszul},\mathrm{LS}}=N^{\mathrm{LS}}(s)^{-1}\cdot\langle\,\cdot\,,\,\cdot\,\rangle^{\mathrm{Koszul}}\)
is defined only on the Eisenstein summand; on the discrete core,
the pairing is unmodified:
\(\langle v,w\rangle^{\mathrm{Koszul},\mathrm{LS}}=\langle v,w\rangle^{\mathrm{Koszul}}\)
for \(v,w\in\mathcal{V}^{\mathrm{prim}}_{\mathrm{disc}}\).
Combining with \Cref{v4-arith:w5-c:thm:discrete-core-identity}
gives (\ref{v4-arith:w6-c:eq:discrete-preserved}).
\end{proof}

\section{Full \(\mathcal{V}^{\mathrm{prim}}\) Koszul--Petersson Identity (Including Eisenstein)}
\label{v4-arith:w6-c:LS:main}

The three compatibilities give the full-sector identity.

\begin{theorem}[full Koszul--Petersson compatibility via LS normalisation]
\label{v4-arith:w6-c:thm:main}
\ClaimStatusProvedHereConditional \emph{(conditional on Ramanujan for Maass
cusp forms on the Maass cuspidal discrete sector; conditional on the
LS Plancherel-measure identification on the Eisenstein continuous
sector)}. For all
\(v,w\in\mathcal{V}^{\mathrm{prim}}\),
\begin{equation}
\label{v4-arith:w6-c:eq:main}
\langle v,w\rangle^{\mathrm{Pet}}\;=\;\langle v,\sigma^{\mathrm{BC}}(w)\rangle^{\mathrm{Koszul},\mathrm{LS}},
\end{equation}
with \(\langle\,\cdot\,,\,\cdot\,\rangle^{\mathrm{Koszul},\mathrm{LS}}\)
the Langlands--Shahidi-corrected Koszul pairing of
\Cref{v4-arith:w6-c:def:LS-Koszul} on the Eisenstein summand and
the unmodified Koszul pairing on the discrete core. The identity
extends
\Cref{v4-arith:w5-c:thm:discrete-core-identity} from the discrete
core to the full \(\mathcal{V}^{\mathrm{prim}}\) including the
Eisenstein continuum under the stated LS normalisation hypothesis.
\end{theorem}

\begin{proof}
Decompose
\(\mathcal{V}^{\mathrm{prim}}=\mathcal{V}^{\mathrm{prim}}_{\mathrm{disc}}\oplus\mathcal{V}^{\mathrm{prim}}_{\mathrm{cont}}\)
by Langlands orthogonality
(Langlands 1966~\S6). On the
discrete core,
\Cref{v4-arith:w6-c:prop:discrete-fixation} gives
(\ref{v4-arith:w6-c:eq:main}) with the conditional on Ramanujan
for Maass inherited from
\Cref{v4-arith:w5-c:thm:discrete-core-identity}. On the
Eisenstein continuum,
\Cref{v4-arith:w6-c:cor:unconditional} supplies the conditional spectral
density normalisation match, which by the Plancherel formula for
the continuous spectrum
(Moeglin--Waldspurger~\S IV.2) yields
(\ref{v4-arith:w6-c:eq:main}) on \(\mathcal{V}^{\mathrm{prim}}_{\mathrm{cont}}\).
Compatibility of the two on their intersection (empty, by
orthogonality) is trivial.
\end{proof}

\begin{corollary}[Eisenstein density mismatch reduced by LS normalisation]
\label{v4-arith:w6-c:cor:mismatch-closed}
\ClaimStatusProvedHereConditional \emph{(conditional on the LS
Plancherel-measure identification)}. The Eisenstein density mismatch of
\Cref{v4-arith:w5-c:prop:mismatch} and
\Cref{v4-arith:kms-gns:cor:eisenstein-final} is reduced under the LS redefinition
\Cref{v4-arith:w6-c:def:LS-Koszul}: the density mismatch on
\(\mathrm{Re}(s)=1/2\) between Koszul and Petersson vanishes
pointwise once the LS-corrected quadratic form is identified with the
positive Plancherel measure. No Ramanujan input is used on the
continuum, but the measure identification is an extra hypothesis.
\end{corollary}

\begin{proof}
\Cref{v4-arith:w6-c:cor:unconditional} gives the pointwise density
match. The domain of the modification is only the Eisenstein
summand
(\Cref{v4-arith:w6-c:rem:discrete-untouched}); the Tate
factorisation (\Cref{v4-arith:w6-c:prop:tate-factorisation}) and
the Heisenberg self-duality
(\Cref{v4-arith:w6-c:prop:koszul-selfduality}) confirm the
redefinition is compatible with the rest of the arithmetic Koszul
structure.
\end{proof}

\section{Source Disjointness and Mathematical-Form Comparison}
\label{v4-arith:w6-c:LS:disjointness-comparison}

\begin{remark}[Source disjointness, conventions, and hypotheses for \Cref{v4-arith:w6-c:thm:main}]
\label{v4-arith:w6-c:rem:disjointness-comparison}
\begin{description}
\item[Source disjointness.] Derivation of the Koszul side:
Beilinson--Drinfeld 2004, Positselski 2009, Bost--Connes 1995.
Verification of the Petersson side: Langlands 1966,
Moeglin--Waldspurger 1995, Shahidi 2010, Casselman--Shalika 1980,
Jacquet--Piatetski-Shapiro--Shalika 1983, Tate 1950, Weil 1967,
Titchmarsh 1986. No source appears on both.
\item[Sign and convention.] The factor
\(N^{\mathrm{LS}}(s)^{-1}=\xi(2s)/\xi(2s-1)\) is a unitary scattering
normalisation on the critical line. The positive measure is supplied
only after the Shahidi Plancherel transport named in
\Cref{v4-arith:w6-c:thm:density-match}. Convention consistent with
Shahidi 2010 Thm~3.5.1 and Moeglin--Waldspurger~\S IV.3.10.
\item[Ambient category.] Both sides of
(\ref{v4-arith:w6-c:eq:main}) are sesquilinear forms on
\(\mathcal{V}^{\mathrm{prim}}\otimes\overline{\mathcal{V}^{\mathrm{prim}}}\)
with values in~\(\mathbb{C}\).
\item[Hypotheses.] The identity is conditional on the LS
Plancherel-measure identification on the Eisenstein continuum and on
Ramanujan for Maass on the Maass cuspidal discrete sector (inherited
from \Cref{v4-arith:w5-c:thm:discrete-core-identity}). Both
conditions are explicit in the theorem statement. \emph{The conditions are explicit.}
\item[Orthogonality.] The Eisenstein summand is
orthogonal to the discrete core under both Petersson (Langlands
1966~\S6) and Koszul (\Cref{v4-arith:w6-c:prop:discrete-fixation}).
The LS correction is supported on the Eisenstein summand only.
\item[Cited results.] The Shahidi 2010 Thm~3.5.1
global normalisation, the Casselman--Shalika 1980 local formula,
and the Moeglin--Waldspurger~\S IV.3.10 unitarity on the critical
line are classical theorems. They do not by themselves identify the
bar--cobar Eisenstein pairing with the positive Plancherel measure.
\emph{With the stated residual.}
\item[Numerical comparison.] No numerical comparison is involved.
\end{description}
Six independent checks yield explicit conditions: Ramanujan for Maass
on the cuspidal discrete sector and the LS Plancherel-measure
identification on the Eisenstein sector.
\end{remark}

\begin{remark}[discrete-core and Eisenstein normalisations]
\label{v4-arith:w6-c:rem:form}
Chapter~\ref{v4-arith:w5-c:full-vprim-P3} (the restricted-class comparison) reduces the
Koszul--Petersson identity on the discrete core and excludes the
Eisenstein continuum by domain restriction. The
Li-criterion path (\Cref{v4-arith:w5-c:cor:Li-path-closes}) uses
this restriction only after the H-P\(^{\mathrm{Ar}}\) determinant-trace
realisation.
The bounded-interval reduction gives a conditional extension to the continuum:
one must insert the LS scalar \(N^{\mathrm{LS}}\) and identify
the resulting quadratic form with the positive Plancherel measure.
\end{remark}

\section{Sources}
\label{v4-arith:w6-c:LS:sources}

The primary sources enter in the following order.

\begin{description}
\item[Shahidi, F.] \emph{Eisenstein Series and Automorphic
\(L\)-Functions}, AMS Colloquium Publications~58, 2010. The
Langlands--Shahidi method for normalising intertwining operators;
Theorem~3.5.1 for the global normalisation we denote
\(N^{\mathrm{LS}}(s)\).
\item[Moeglin, C. and Waldspurger, J.-L.]
\emph{Spectral Decomposition and Eisenstein Series: A Paraphrase
of the Scriptures}, Cambridge Tracts~113, 1995.
\S II.1.7 for the intertwining operator \(M(s)\);
\S IV.2 for the Plancherel formula;
\S IV.3.10 for unitarity of the normalised intertwiner on the
critical line.
\item[Langlands, R. P.] \emph{On the Functional Equations
Satisfied by Eisenstein Series}, Lecture Notes in Math.~544,
Springer 1976 (manuscript 1966). Foundational spectral
decomposition of \(L^{2}([GL_{n}])\) into discrete and continuous
summands; Maass--Selberg relations; functional equation of
Eisenstein series.
\item[Casselman, W. and Shalika, J.] \emph{The unramified principal
series of \(p\)-adic groups. II. The Whittaker function},
Compositio Math.~41 (1980), 207--231. Unramified spherical
Whittaker formula used for the local factor
\(N^{\mathrm{LS}}_{p}(s)=\zeta_{p}(2s-1)/\zeta_{p}(2s)\) at
unramified finite~\(p\).
\item[Jacquet, H., Piatetski-Shapiro, I., and Shalika, J.]
\emph{Rankin--Selberg convolutions}, Amer.~J.~Math.~105 (1983),
367--464. Local \(\varepsilon\)-factors at ramified principal series.
\item[Tate, J.] \emph{Fourier analysis in number fields and Hecke's
zeta functions}, Ph.D.~thesis, Princeton 1950;
reprinted Cassels--Fr\"ohlich 1967 Ch.~XV. Adelic restricted
tensor product; local-global factorisation of completed zeta.
\item[Weil, A.] \emph{Basic Number Theory}, Springer 1967.
Tamagawa normalisation.
\item[Titchmarsh, E. C.] \emph{The Theory of the Riemann
Zeta-Function} (2nd ed., revised D.~R.~Heath-Brown), Oxford
University Press 1986. Functional equation of~\(\xi\); Euler
product for~\(\zeta\).
\item[Bost, J.-B. and Connes, A.] \emph{Hecke algebras, type III
factors and phase transitions with spontaneous symmetry breaking
in number theory}, Selecta Math.~1 (1995), 411--457. The
Bost--Connes \(C^{*}\)-algebra; KMS\({}_{1}\) state; involution
\(\sigma^{\mathrm{BC}}\).
\item[Beilinson, A. and Drinfeld, V.] \emph{Chiral Algebras}, AMS
Colloquium~51, 2004. Chiral bar--cobar construction and
parabolic-induction compatibility (\S3.5).
\item[Positselski, L.] \emph{Two kinds of derived categories,
Koszul duality, and comodule-contramodule correspondence},
arXiv:0905.2621, 2009. Koszul duality input for the arithmetic
bar--cobar.
\end{description}

\paragraph{Disjoint sources.}
The proof sources for the arithmetic Koszul pairing
(Beilinson--Drinfeld 2004, Positselski 2009, Bost--Connes 1995)
are disjoint from the analytic sources for the
Langlands--Shahidi normalisation and Plancherel measure
(Shahidi 2010, Moeglin--Waldspurger 1995, Langlands 1966,
Casselman--Shalika 1980, Jacquet--Piatetski-Shapiro--Shalika
1983, Tate 1950, Weil 1967, Titchmarsh 1986). No source appears
on both sides of the identification.
